% !TEX root = task-6-2-1.tex
% Editable task bank. Append new BankTask blocks at the end of the document.
\documentclass[12pt,a4paper]{article}
\usepackage{physics-textbook}
\usepackage{graphicx,needspace}
\usetikzlibrary{decorations.pathmorphing}
\newif\ifShowSolutions
\ifdefined\TutorEdition
  \ShowSolutionstrue
\else
  \ShowSolutionsfalse % Change to \ShowSolutionstrue to display solutions.
\fi
\newcounter{banktask}
\newsavebox{\BankTaskBox}
\newcommand{\BankTaskContent}[3]{%
  \noindent\textbf{\thebanktask.}\quad #2\par
  \ifShowSolutions
    \Needspace{50mm}% Keep the solution heading with its opening text or diagram.
    #3\par
    \ifdefined\TutorEdition
      \ifcsname TutorSource#1\endcsname
        \nopagebreak\smallskip{\small\csname TutorSource#1\endcsname\par}%
      \fi
    \fi
  \fi
}
\newcommand{\BankTask}[3]{%
  \refstepcounter{banktask}\par\medskip
  \sbox{\BankTaskBox}{\begin{minipage}{\linewidth}%
    \BankTaskContent{#1}{#2}{#3}\end{minipage}}%
  \ifdim\dimexpr\ht\BankTaskBox+\dp\BankTaskBox\relax>\dimexpr\textheight-20mm\relax
    \Needspace{35mm}\BankTaskContent{#1}{#2}{#3}%
  \else
    \Needspace{\dimexpr\ht\BankTaskBox+\dp\BankTaskBox\relax}%
    \noindent\usebox{\BankTaskBox}\par
  \fi
  \medskip
}
\begin{document}
\BookSubsection{6.2.1}{Энергия фотона}


% Source sdamgia-2003: https://phys-ege.sdamgia.ru/problem?id=2003
\BankTask{sdamgia-2003}{
Длина волны рентгеновского излучения равна $10 ^{ - 10 } \,\text{м}.$ Во сколько раз энергия одного фотона этого излучения превосходит энергию фотона видимого света длиной волны $4 \cdot 10 ^{ - 7 } \,\text{м}$?
}{
\smallskip\noindent\textbf{Решение}\par
Энергия фотона: $E=h\nu.$ Учитывая связь $\lambda \nu =c,$ получаем выражение для энергии фотона через длину волны: $E= \frac{hc}{\lambda} .$ Таким образом, энергия рентгеновского фотона больше энергии фотона видимого света в 

\[\frac{E_{\,\text{рент}}.}{E_{\,\text{вид}}. \,\text{св}.} = \frac{\lambda _{\,\text{вид}}. \,\text{св}.}{\lambda _{\,\text{рент}}.} = \frac{4 \cdot 10 ^{ - 7 }}{10 ^{ - 10 }} =4000\,\text{раз}.\]

Ответ: $4000$.
}

% Source sdamgia-2016: https://phys-ege.sdamgia.ru/problem?id=2016
\BankTask{sdamgia-2016}{
Атом испустил фотон с энергией $6 \cdot 10 ^{ - 18 } \,\text{Дж}.$ На сколько уменьшится импульс атома? (в $10 ^{ - 26 } \,\text{кг} \cdot \,\text{м/с}.$)
}{
\smallskip\noindent\textbf{Решение}\par
Энергия и импульс фотона связаны соотношением $E=pc.$ Следовательно, фотон унес импульс

\[p= \frac{E}{c} = \frac{6 \cdot 10 ^{ - 18 } \,\text{Дж}}{3 \cdot 10 ^{8} \,\text{м/с}} =2 \cdot 10 ^{ - 26 } \,\text{кг} \cdot \,\text{м/с}.\]

Согласно закону сохранения импульса, полный импульс системы не должен изменяться. Следовательно, импульс атома уменьшился на $2 \cdot 10 ^{ - 26 } \,\text{кг} \cdot \,\text{м/с}.$

Ответ: $2$.
}

% Source sdamgia-3196: https://phys-ege.sdamgia.ru/problem?id=3196
\BankTask{sdamgia-3196}{
Фотон с энергией $E$ движется в вакууме. Пусть $h$ — постоянная Планка, $c$ — скорость света в вакууме. Чему равны частота и импульс фотона?
\TaskMatchingColumns{Физическая величина}{А) Частота фотона\\Б) Импульс фотона}{Выражение для неё}{1) $hc/E$\\2) $E/c^2$\\3) $E/c$\\4) $E/h$}
\begin{center}\begin{tabular}{|c|c|}\hline А&Б\\\hline &\\\hline\end{tabular}\end{center}
}{
\smallskip\noindent\textbf{Решение}\par
Энергия фотона связана с частотой света соотношением $E=h\nu .$ Следовательно, частота фотона равна $\nu = \frac{E}{h}$ (А  — $4$). Импульс и энергия фотона связаны соотношением $E=pc.$ Таким образом, импульс фотона равен $p= \frac{E}{c}$ (Б  — $3$).

Ответ: $43$.
}

% Source sdamgia-5160: https://phys-ege.sdamgia.ru/problem?id=5160
\BankTask{sdamgia-5160}{
Энергия фотона в рентгеновском дефектоскопе в $2$ раза больше энергии фотона в рентгеновском медицинском аппарате. Каково отношение частоты электромагнитных колебаний в первом пучке рентгеновских лучей к частоте во втором пучке?
}{
\smallskip\noindent\textbf{Решение}\par
Энергия фотона связана с частотой соотношением $E=h\nu.$ Таким образом, искомое отношение частоты электромагнитных колебаний в первом пучке рентгеновских лучей к частоте во втором пучке равно

\[\frac{\nu_1}{\nu_2} = \frac{E_1}{E_2} =2.\]

Ответ: $2$.
}

% Source sdamgia-5195: https://phys-ege.sdamgia.ru/problem?id=5195
\BankTask{sdamgia-5195}{
Модуль импульса фотона в рентгеновском дефектоскопе $2$ раза больше модуля импульса фотона в рентгеновском медицинском аппарате. Каково отношение энергии фотона в первом пучке рентгеновских лучей к энергии фотона во втором пучке?
}{
\smallskip\noindent\textbf{Решение}\par
Энергия фотона связана с импульсом соотношением $E=pc.$ Таким образом, искомое отношение энергии фотона в первом пучке рентгеновских лучей к энергии фотона во втором пучке равно

\[\frac{E_1}{E_2} = \frac{p_1}{p_2} =2.\]

Ответ: $2$.
}

% Source sdamgia-7361: https://phys-ege.sdamgia.ru/problem?id=7361
\BankTask{sdamgia-7361}{
В вакууме распространяются два параллельных пучка света. Свет первого пучка характеризуется длиной волны $300\,\text{нм}$, а свет второго пучка  — частотой $0{,}5 \cdot 10 ^{ 15 } \,\text{Гц}.$ Во сколько раз отличается энергия фотона из первого пучка от энергии фотона из второго пучка? Скорость света принять равной $3$·$10$$^{8}$ м/с.
}{
\smallskip\noindent\textbf{Решение}\par
Энергия фотона равна $E=h\nu= \frac{hc}{\lambda} .$ Отношение $\frac{E_1}{E_2} = \frac{\frac{hc}{\lambda_1}}{h\nu_2} = \frac{c}{\lambda_1\nu_2} = \frac{3 \cdot 10 ^{8}}{300 \cdot 10 ^{ - 9 } \cdot 0{,}5 \cdot 10 ^{ 15 }} =2.$

Ответ: $2$.
}

% Source sdamgia-7866: https://phys-ege.sdamgia.ru/problem?id=7866
\BankTask{sdamgia-7866}{
Зеленый свет ($\lambda$ = $550\,\text{нм}$) переходит из воздуха в стекло с показателем преломления $1{,}5$. Определите отношение энергии фотона в воздухе к его энергии в стекле.
}{
\smallskip\noindent\textbf{Решение}\par
Энергия фотона равна $E=h\nu.$ При переходе из одной среды в другую частота фотона $\nu$ не меняется, значит, и энергия остается неизменной. Отношение энергий равно единице.

Ответ: $1$.
}

% Source sdamgia-6221: https://phys-ege.sdamgia.ru/problem?id=6221
\BankTask{sdamgia-6221}{
Мощность излучения лазерной указки с длиной волны $\lambda$  =  $600\,\text{нм}$ равна P  =  $2$ мВт. Определите число фотонов, излучаемых указкой за $1\,\text{с}$.
}{
\smallskip\noindent\textbf{Решение}\par
Один фотон света с частотой $\nu$ обладает энергией $h\nu=h \frac{c}{\lambda} .$ Энергия излучаемая за время t указкой  — $E=Pt.$ Значит, число фотонов N, излучаемых указкой за время $t:$

\[N= \frac{E}{h\nu} = \frac{P\lambda t}{hc} = \frac{2 \cdot 10 ^{ - 3 } \cdot 1 \cdot 600 \cdot 10 ^{ - 9 }}{6{,}6 \cdot 10 ^{ - 34 } \cdot 3 \cdot 10 ^{8}} \approx0{,}6 \cdot 10 ^{ 16 } =6 \cdot 10 ^{ 15 } \,\text{фотонов}.\]

Ответ: $6 \cdot 10 ^{ 15 } \,\text{фотонов}.$
}

% Source sdamgia-29865: https://phys-ege.sdamgia.ru/problem?id=29865
\BankTask{sdamgia-29865}{
Какой должна быть длина волны фотона, чтобы его энергия равнялась кинетической энергии протона, движущегося со скоростью $2$ · $10$$^{4}$ м/c?
}{
\smallskip\noindent\textbf{Решение}\par
Кинетическая энергия протона равна $E_k= \frac{m_p v ^2}{2} .$ Энергия фотона определяется формулой $E_{\,\text{ф}}=h v = \frac{h c}{\lambda} .$ Учитывая, что эти энергии равны, получаем $\frac{m_p v ^2}{2} = \frac{h c}{\lambda} .$ Откуда 

\[\lambda= \frac{2 h c}{m_p v ^2} = \frac{2 \cdot 6{,}6 \cdot 10 ^{ - 34 } \cdot 3 \cdot 10 ^{ 8 }}{1{,}673 \cdot 10 ^{ - 27 } \cdot 4 \cdot 10 ^{ 8 }} \approx5{,}92 \cdot 10 ^{ - 7 } \,\text{м}=592 \,\text{нм}.\]

Ответ: $592\,\text{нм}$.
}

% Source sdamgia-3723: https://phys-ege.sdamgia.ru/problem?id=3723
\BankTask{sdamgia-3723}{
Энергия ионизации атома кислорода равна $14\,\text{эВ}$. Найдите максимальную длину волны света, которая может вызвать ионизацию атома кислорода.

Справочные данные: постоянная Планка $h=6{,}6 \cdot 10 ^{ - 34 } \,\text{Дж} \cdot \,\text{с},$ $1\,\text{эВ}=1{,}6 \cdot 10 ^{ - 19 } \,\text{Дж}.$ (В нанометрах; округлите до целого.)
}{
\smallskip\noindent\textbf{Решение}\par
Длина волны связана с частотой и скоростью света соотношением: $\lambda\nu=c.$ Следовательно, максимально возможной длине волны, соответствует минимально возможная частота. 
Согласно постулатам Бора, для перехода электрона на более высокий уровень необходимо, чтобы атом поглотил квант энергии, равный по величине разности энергий конечного и начального состояний. Для ионизации атома необходимо, чтобы энергия поглощаемого фотона была не меньше энергии ионизации: $E_{\,\text{фот}} > \,\text{или} \,\text{равно} E_{\,\text{ион}}.$ Энергия фотона пропорциональна частоте света: $E_{\,\text{фот}}=h\nu.$ Таким образом, максимальная длина волны света, которая может вызвать ионизацию атома кислорода равна: 

\[\lambda= \frac{c}{\nu} = \frac{\,\text{с}}{E_{\,\text{ион}}/h} = \frac{hc}{E_{\,\text{ион}}} = \frac{6{,}6 \cdot 10 ^{ - 34 } \,\text{Дж} \cdot \,\text{с} \cdot 3 \cdot 10 ^{8} \,\text{м/с}}{14 \cdot 1{,}6 \cdot 10 ^{ - 19 } \,\text{Дж}} \approx 0{,}88 \cdot 10 ^{ - 7 } \,\text{м}=88\,\text{нм}.\]

Ответ: $88\,\text{нм}$.
}


% Source 7D7448: https://ege.fipi.ru/bank/questions.php?search=1&proj=BA1F39653304A5B041B656915DC36B38&qid=7D7448
\BankTask{7D7448}{
Установите соответствие между физическими величинами и формулами, по которым их можно рассчитать ($E$ — энергия фотона; $h$ — постоянная Планка; $p$ — импульс фотона).

\TaskMatchingColumns[.53\linewidth]{Физические величины}{
А) длина волны фотона\par
Б) частота фотона
}{Формулы}{
1) $\dfrac{p}{h}$\par
2) $\dfrac{h}{p}$\par
3) $\dfrac{E}{p}$\par
4) $\dfrac{E}{h}$
}

\medskip
\begin{tabular}{|c|c|}
\hline
А & Б \\
\hline
\phantom{2} & \phantom{4} \\
\hline
\end{tabular}
}{
\smallskip\noindent\textbf{Решение от GPT}\par
Из соотношения $p=h/\lambda$ получаем $\lambda=h/p$, поэтому А соответствует формуле 2. По формуле Планка $E=h\nu$, откуда $\nu=E/h$, поэтому Б соответствует формуле 4. Ответ: $24$.
}

% Source 64FC99: https://ege.fipi.ru/bank/questions.php?search=1&proj=BA1F39653304A5B041B656915DC36B38&qid=64FC99
\BankTask{64FC99}{
Установите соответствие между физическими величинами и формулами, по которым их можно рассчитать ($\lambda$ — длина волны фотона; $c$ — скорость света в вакууме; $h$ — постоянная Планка).

\TaskMatchingColumns[.53\linewidth]{Физические величины}{
А) частота фотона\par
Б) импульс фотона
}{Формулы}{
1) $\dfrac{h}{\lambda}$\par
2) $\dfrac{h\lambda}{c}$\par
3) $\dfrac{c}{\lambda}$\par
4) $c\lambda$
}

\medskip
\begin{tabular}{|c|c|}
\hline
А & Б \\
\hline
\phantom{3} & \phantom{1} \\
\hline
\end{tabular}
}{
\smallskip\noindent\textbf{Решение от GPT}\par
Из $c=\lambda\nu$ следует $\nu=c/\lambda$, поэтому А соответствует формуле 3. Импульс фотона $p=h/\lambda$, поэтому Б соответствует формуле 1. Ответ: $31$.
}

% Source D4FFF9: https://ege.fipi.ru/bank/questions.php?search=1&proj=BA1F39653304A5B041B656915DC36B38&qid=D4FFF9
\BankTask{D4FFF9}{
Установите соответствие между графиками, представленными на рисунках, и зависимостями, которые они могут выражать.

\par\medskip\TaskMatchHeading{Графики}

А) рисунок 1

Б) рисунок 2

\par\medskip\TaskMatchHeading{Зависимости}

1) зависимость силы фототока от напряжения между электродами при неизменной освещённости

2) зависимость энергии фотона от длины волны

3) зависимость максимальной энергии фотоэлектронов от частоты света

4) зависимость энергии фотона от частоты света

\begin{tabular}{|c|c|}
\hline
А & Б \\
\hline
\phantom{2} & \phantom{3} \\
\hline
\end{tabular}

\medskip
\begin{center}
\begin{tikzpicture}[book diagram, x=0.024cm, y=0.024cm]
\draw[-{book tip}] (0,0) -- (150,0);
\draw[-{book tip}] (0,0) -- (0,127);
\node[below left] at (0,0) {$0$};
\draw[line width=0.5pt] (21,125) .. controls (22,118) and (21,105) .. (23,94) .. controls (26,73) and (32,58) .. (42,48) .. controls (54,36) and (69,29) .. (87,25) .. controls (104,21) and (125,19) .. (143,19);
\draw (0,103) -- (22,103);
\draw (22,103) -- (22,80);
\draw (22,32) -- (22,8);
\draw (0,23) -- (7,23);
\draw (32,23) -- (55,23);
\draw (80,23) -- (104,23);
\draw (104,23) -- (104,0);
\node[below] at (72,-18) {А) Рисунок 1};
\end{tikzpicture}
\hspace{1cm}
\begin{tikzpicture}[book diagram, x=0.024cm, y=0.024cm]
\draw[-{book tip}] (0,0) -- (163,0);
\draw[-{book tip}] (0,0) -- (0,127);
\node[below left] at (0,0) {$0$};
\draw[line width=0.7pt] (0,0) -- (52,0) -- (151,99);
\node[below] at (76,-18) {Б) Рисунок 2};
\end{tikzpicture}
\end{center}
}{
\smallskip\noindent\textbf{Решение от GPT}\par
График А — убывающая гипербола: $E=hc/\lambda$, поэтому ему соответствует зависимость 2. График Б описывает линейный рост максимальной энергии фотоэлектронов выше пороговой частоты: $K_{\max}=h\nu-A_{\text{вых}}$, то есть зависимость 3. Зависимость 1 представляет рост фототока с выходом на насыщение, а зависимость 4 — прямую $E=h\nu$, проходящую через начало координат; они не соответствуют данным графикам. Ответ: $23$.
}

% Source 369826: https://ege.fipi.ru/bank/questions.php?search=1&proj=BA1F39653304A5B041B656915DC36B38&qid=369826
\BankTask{369826}{
На металлическую пластинку падает пучок монохроматического света. При этом наблюдается явление фотоэффекта.

На графике А представлена зависимость энергии фотонов, падающих на катод, от физической величины \(x_1\), а на графике Б — зависимость максимальной кинетической энергии фотоэлектронов от физической величины \(x_2\).

Какая из физических величин отложена на горизонтальной оси на графике А и какая — на графике Б?

\TaskMatchingColumns[.53\linewidth]{Графики}{
А) график А\par
Б) график Б
}{Физическая величина \(x\)}{
1) частота\par
2) длина волны\par
3) массовое число\par
4) заряд ядра
}

\medskip
\begin{tabular}{|c|c|}
\hline
А & Б \\
\hline
\phantom{2} & \phantom{1} \\
\hline
\end{tabular}

\medskip
\begin{tikzpicture}[book diagram,x=0.45cm,y=0.23cm]
\node at (3.5,20.5) {А)};
\draw[-{book tip}] (0,0) -- (7.3,0) node[below] {$x_1$};
\draw[-{book tip}] (0,0) -- (0,18.5) node[left] {$E$};
\node[below left] at (0,0) {$0$};
\foreach \y in {0.5,1,...,17.5} {\draw (-0.09,\y) -- (0.09,\y);}
\foreach \y in {5,10,15} {\node[left] at (-0.14,\y) {$\y$};}
\foreach \xx/\yy in {1/15,2/7.5,3/5,4/3.75,5/3,6/2.5} {
\draw[thin] (0,\yy) -- (\xx,\yy) -- (\xx,0);
}
\draw[thick] plot[domain=0.82:7,samples=100] (\x,{15/\x});
\foreach \x/\y in {1/15,2/7.5,3/5,4/3.75,5/3,6/2.5}{\fill (\x,\y) circle[radius=1pt];}
\end{tikzpicture}
\qquad
\begin{tikzpicture}[book diagram,x=0.45cm,y=0.45cm]
\node at (3.5,5.25) {Б)};
\draw[-{book tip}] (0,0) -- (7.1,0) node[below] {$x_2$};
\draw[-{book tip}] (0,-1.25) -- (0,4.7) node[left] {$E$};
\node[left] at (0,0) {$0$};
\draw[thick] (0,0) -- (2,0) -- (5.7,3.7);
\end{tikzpicture}
}{
\smallskip\noindent\textbf{Решение от GPT}\par
Энергия фотона \(E=h\nu=hc/\lambda\), поэтому убывающая гипербола А соответствует длине волны (2). По уравнению Эйнштейна \(K_{\max}=h\nu-A_{\text{вых}}\), поэтому возрастающая прямая Б соответствует частоте (1). Массовое число (3) и заряд ядра (4) не являются переменными в этих зависимостях. Ответ: А — 2, Б — 1.
}

\end{document}
